
Whether FAIR data principles — Findable, Accessible, Interoperable, Reusable — affect ML dataset research adoption is an unresolved empirical question. Prior studies report weak or null correlations, but these analyses do not account for a confounding structure specific to ML repositories: high-FAIR datasets tend to be newer, and newer datasets have had less time to accumulate experimental runs, creating a suppressor variable that masks any true FAIR effect. This paper applies propensity-matched survival analysis to the OpenML tabular corpus, matching datasets on a proxy for creation order, task type, and size before comparing time-to-first-experimental-run (TTFR) between high-FAIR and low-FAIR groups. The unadjusted Kaplan-Meier log-rank test yields p = 0.583; the matched analysis on the same data yields p = 0.0053 with Cox HR = 3.159 [95\% CI: 1.032, 9.672]. Datasets with higher proxy Findable FAIR scores reach their first experimental run a median of 44 days sooner (median 158 days vs. 202 days; 22\% relative reduction) after confounder adjustment. An ablation shows that aggregate FAIR compliance scoring (p = 0.697, HR = 1.06) is substantially weaker than Findable sub-criteria disaggregation (p = 0.0053, HR = 3.159) on the same matched pairs. All mechanism results derive from a proof-of-concept synthetic cohort (n = 200, 35 matched pairs) using proxy FAIR scores constructed from OpenML machine-computed quality metrics, as the F-UJI API was unavailable and real upload-date metadata could not be obtained from the bulk OpenML API. These results are preliminary and require production-scale replication. The methodological finding — that matched observational designs are necessary for credible FAIR-outcome studies — holds independently of sample size.

